% ---------------------------------------------------------------
% Replaces the final paragraph of
% "Literature on relevant finite discretization methods".
% VERSION A -- solution only. Safe with the results you have reported.
% ---------------------------------------------------------------

In this work we bootstrap the level-set based finite volume method on Cartesian grids
proposed by Bochkov \& Gibou (2020)~\cite{bochkov2020}. The method rests on Taylor
expansions in the normal direction to relate the solution at a grid point to its value
on the interface, combined with geometric integration over the cut cells the interface
crosses. Applied to jump conditions, it yields second-order accurate solutions with
first-order accurate gradients in the $L^{\infty}$-norm. In the Robin setting treated
here the same construction is used to eliminate the unknown interface value: writing
the boundary condition along the normal and taking one Taylor step from the cell centre
gives
\begin{equation}
    u_{\Gamma} = \frac{\mu\, u_{ijk} + |\delta_{ijk}|\, g}
                      {\mu + \alpha\, |\delta_{ijk}|},
    \label{eq:taylor_robin}
\end{equation}
where $\delta_{ijk}$ is the signed distance from $\mathbf{x}_{ijk}$ to $\Gamma$. This
projection is first-order accurate in $\delta_{ijk}$, and since it is the sole
mechanism by which the boundary condition enters the discretisation, it limits the
scheme to first order overall irrespective of the accuracy of the interior stencil.
Our numerical results in Section~\ref{sec:results} are consistent with this: we
observe first-order convergence of the solution in both the root-mean-squared and
$L^{\infty}$ norms across all three geometries.


% ---------------------------------------------------------------
% VERSION B -- use INSTEAD of the last sentence above, once you have
% extracted gradient errors from the run logs and computed their orders.
% Replace ORDER_HERE with the measured value.
% ---------------------------------------------------------------
%
% Our numerical results in Section~\ref{sec:results} are consistent with this: we
% observe first-order convergence of the solution in both the root-mean-squared and
% $L^{\infty}$ norms across all three geometries, together with ORDER_HERE-order
% convergence of the solution gradient in the $L^{\infty}$-norm.
